\bookchapter{פרק 17: קישוריות הטלפון עם מכשירים אחרים}

\chapterintro{להגיד את האמת, מה שהכי הטריד אותי פעם זה שצילמתי תמונה בטלפון, רציתי להעביר למחשב לעריכה, והייתי צריך למצוא כבל, לחכות להעברה. אחר כך גיליתי שבכלל לא צריך כבל, קישוריות טובה מהירה פי עשרה מכבל.}

\booksection{אקוסיסטם Apple (AirDrop + לוח גזירה אוניברסלי)}

: אם אתה משתמש ב-iPhone עם Mac, AirDrop היא הפונקציה הראשונה שכדאי ללמוד. בטלפון לחץ על שיתוף AirDrop בחר את ה-Mac שלך, התמונה עוברת מיד, בלי תעבורה בלי Wi-Fi, משתמש ב-Bluetooth ו-Wi-Fi Direct. יש גם פונקציה שקל לפספס בשם “לוח גזירה אוניברסלי” - מעתיק טקסט בטלפון, ב-Mac פשוט מדביקים וזה מופיע, כמו קסם.

\begin{itemize}
  \item \textbf{אנדרואיד + Windows (Link to Windows)}: ל-Microsoft יש אפליקציה בשם "Link to Windows", מובנית בטלפוני Samsung, לשאר טלפוני אנדרואיד יורדים מהחנות. אחרי חיבור אפשר לראות במחשב התראות טלפון, לענות ל-SMS, לדפדף בתמונות, אפילו לקבל שיחות. אני משתמש כבר שנתיים, בפגישות לא צריך להסתכל על הטלפון כל הזמן, התראה על המסך מספיקה.
  \item \textbf{קישוריות מולטי-מסך מותגית}: "Multi-Screen Collaboration" של Huawei ו-"Xiaomi Smart Connect" של Xiaomi יותר קיצוניים - משליכים את מסך הטלפון ישירות למחשב, אפשר לגרור קבצים בעכבר מהטלפון למחשב, וגם לגלול ב-Douyin או לענות ל-WeChat במחשב. העתקה-הדבקה חוצת מכשירים, גרירת קבצים - חוויה חלקה מאוד.
  \item \textbf{שתיים, הקרנה לטלוויזיה: כל המשפחה צופה יחד}בחזרה הביתה לראש השנה, בת דודתי רצתה להקרין את סרטון טיול הסיום שלה בטלוויזיה. אחרי התעסקות ארוכה, התברר שזה פשוט מאוד:
\end{itemize}

\booksection{iPhone טלוויזיה}

: משך כלפי מטה מהפינה הימנית העליונה את מרכז הבקרה, לחץ על הסמל של שני מלבנים מונחים אחד על השני (שיקוף מסך), מצא את שם הטלוויזיה שלך, לחיצה אחת ומחובר. התנאי: הטלוויזיה והטלפון מחוברים לאותו Wi-Fi, והטלוויזיה תומכת ב-AirPlay (רוב הטלוויזיות החכמות תומכות).

\begin{itemize}
  \item \textbf{אנדרואיד טלוויזיה}: כל מותג קורא לזה אחרת, Huawei קורא "הקרנה אלחוטית", Xiaomi קורא "הקרנה", OPPO/vivo קוראים "אינטראקציה מולטי-מסך". בדרך כלל ניתן למצוא בפאנל המתגים המהירים, אם לא מוצאים אז לחפש "הקרנה" בהגדרות.
  \item \textbf{הקרנה ישירות מאפליקציות וידאו}: אם רק רוצה לראות בטלוויזיה את iQiyi, Bilibili, YouTube, זה אפילו יותר פשוט - בממשק הנגן חפש סמל קטן של מלבן עם סמל Wi-Fi, לחץ, בחר טלוויזיה, הוידאו עובר. אז הטלפון יכול לעשות דברים אחרים, הוידאו ממשיך בטלוויזיה, ללא הפרעה הדדית.
  \item \textbf{מה עושים אם ההקרנה לא מצליחה}: 90\% מהמקרים זה כי לא מחוברים לאותו Wi-Fi. עוד 10\% זה שפונקציית ההקרנה של הטלוויזיה לא מופעלת, חפש "אינטראקציה מולטי-מסך" או "AirPlay" בהגדרות הטלוויזיה.
  \item \textbf{שלוש, עם שעון/צמיד חכם: מנהל בריאות על שורש כף היד}קניתי לאמי צמיד חכם, בהתחלה היא אמרה "אני לא עושה ספורט, למה לקנות את זה". אחרי שבועיים היא גילתה שהיא ישנה רק חמש שעות וקצת ביום, והתחילה להתאים את שגרת השינה באופן עצמאי. זהו הערך הפרקטי ביותר של שעון/צמיד - להפוך נתוני בריאות בלתי נראים למספרים גלויים.
\end{itemize}

\booksection{מסנן התראות}

: כשהטלפון רועד בלי הפסקה בתיק או בכיס, הרמת שורש כף יד להצצה - האם זה WeChat מהבוס או קידום מ-Pinduoduo, החלטה בשנייה, הודעות שלא שוות להוציא את הטלפון - פשוט תניח להן.

\begin{itemize}
  \item \textbf{סנכרון נתוני בריאות}: השעון מנטר 24 שעות דופק, חמצן בדם, שינה, צעדים, הנתונים האלה מצטברים אוטומטית באפליקציית הבריאות בטלפון. עם הזמן תראה את המגמות שלך - האם הלכת פחות מדי החודש, האם איכות השינה השתפרה השבוע.
  \item \textbf{סנכרון כרטיס תחבורה/כרטיס גישה}: כרטיסי אוטובוס וכרטיסי גישה שהוגדרו בטלפון, רוב המותגים מסנכרנים אוטומטית לשעון. ברכבת התחתית לא צריך להוציא את הטלפון, שורש כף היד קרוב לשער ועברת, באמת יש תחושה של "העתיד כבר כאן".
  \item \textbf{מוצא הטלפון}: לרוב השעונים יש פונקציית "מצא טלפון", לחיצה אחת והטלפון מצלצל. אני משתמש בזה הרבה יותר ממה שציפיתי - שם את הטלפון ומיד שוכח איפה.
  \item \textbf{ארבע, עם בית חכם: הטלפון הוא "המתג הראשי" של הבית}בית חכם נשמע מיסטי, אבל בעצם מדובר בשלושה דברים: שליטה מרחוק בטלפון, הגדרת כללי אוטומציה, שליטה קולית.
\end{itemize}

\booksection{אספר לך על כמה שימושים מעשיים שלי, ותבין כמה זה פרקטי:}

מזגן

: בקיץ בדרך הביתה מפעילים מרחוק את המזגן בטלפון, נכנסים הביתה וקריר. תמורת מאה וקצת יואן קונים שקע מצמד מזגן וזהו, לא צריך להחליף מזגן.

\begin{itemize}
  \item \textbf{אור}: לפני השינה במיטה, עצל לקום לכבות, לוחצים בטלפון או קוראים "היי Siri/Xiao Ai, כבה אור".
  \item \textbf{רובוט שואב}: לפני היציאה לוחצים בטלפון "התחל ניקוי", חוזרים והרצפה נקייה. אפשר גם לקבוע ניקוי יומיומי.
  \item \textbf{מצלמה}: אם יש בבית קשישים, ילדים או חיות מחמד, אפשר להציץ בכל רגע במצב הבית.
  \item \textbf{איזה אקוסיסטם לבחור? בקיצור:}משתמש ב-iPhone העדף מכשירים התומכים ב-Apple HomeKit (קצת יותר יקר אבל פרטיות טובה יותר).
\end{itemize}

הרבה מכשירים, מחפש תמורה Xiaomi Mi Home האקוסיסטם השלם ביותר, יש אור, שקע, שואב וכו'.

\begin{itemize}
  \item משתמש ב-Huawei העדף HarmonyOS Connect, חוויה מיטבית עם מכשירי Huawei.
  \item חמש, אוזניות ורמקולי Bluetooth: חופש אלחוטי
  \item המעבר מאוזניות עם חוט ל-Bluetooth היה אחד השדרוגים המורגשים ביותר בחוויית הטלפון שלי. כבר לא מוציאים כדור חוטים מהתיק, ולא דואגים שהחוט ייתפס בידית הדלת.
\end{itemize}

\booksection{חיבור פשוט}

: אוזניות בהפעלה ראשונה בדרך כלל נכנסות אוטומטית למצב חיבור (נורה מהבהבת), הטלפון פותח Bluetooth ומוצא אותן, לחיצה אחת ומחובר. אחר כך כל פעם שמוציאים - מתחבר אוטומטית, לא צריך לעשות כלום.

\begin{itemize}
  \item \textbf{באותו מותג יש בונוס}: AirPods עם iPhone יציגו אוטומטית חלון עם מצב סוללה, וגם יחליפו בצורה חלקה בין מכשירי Apple. Huawei FreeBuds עם טלפון Huawei, אוזניות Xiaomi עם טלפון Xiaomi - גם יש חלונות ואופטימיזציות דומות.
  \item \textbf{רמקול Bluetooth}: במפגש משפחתי, בקמפינג, רמקול Bluetooth אחד מכפיל את האווירה. אחרי חיבור, הטלפון הוא עמדת ה-DJ שלך.
  \item \textbf{שים לב להשהיה}: במשחקים עם אוזניות Bluetooth עשויה להיות השהיה קלה (חצי שנייה אחרי הירי שומעים את הירי). אם רגיש להשהיה, או אוזניות עם חוט, או קנה אוזניות Bluetooth עם "מצב השהיה נמוכה".
  \item \textbf{שש, עם מערכת רכב: שותף בטוח בנהיגה}להסתכל בטלפון בנהיגה זה לא חוקי, וגם מסוכן לחיים. אבל לפעמים באמת צריך ניווט, לענות לשיחה, להשמיע מוזיקה. כאן נכנסים CarPlay (Apple) ו-Android Auto (אנדרואיד).
\end{itemize}

\begin{tcbinfo}{UWB - "תחושת מרחב" מדויקת מ-Bluetooth}
מלבד Bluetooth, יש שבב UWB. יכולת:\textbf{מדידת מרחק מדויקת ומיקום מרחבי}—מרחק וכיוון, דיוק סנטימטר. Bluetooth “בקרבה”, UWB “1.2 מ' שמאל”. שימושים:

\begin{itemize}
  \item \textbf{AirTag/חיפוש מדויק}:iPhone AirTag "7 מ', שמאל" = UWB. Bluetooth רק "בטווח".
  \item \textbf{מפתח רכב דיגיטלי}:UWB = מפתח רכב, פתיחה/נעילה, מניעת relay.
  \item \textbf{ניווט מקורה}:ניווט "עד פרק" בקניון/שדה, מדויק מ-GPS.
  \item \textbf{דגמים נתמכים}:iPhone 11+ (U1/U2); Xiaomi, Samsung, vivo, OPPO דגלים.
\end{itemize}

\end{tcbinfo}

\begin{tcbinfo}{LE Audio ו-Auracast - דור חדש}
A2DP הוותיק מוחלף על ידי\textbf{LE Audio}(LE Audio), Bluetooth 5.2+:

\begin{itemize}
  \item \textbf{חיסכון בסוללה, איכות צליל טובה יותר}:LC3, איכות טובה מ-SBC, פחות צריכה, יותר סוללה.
  \item \textbf{אודיו דו-ערוצי}: העברה עצמאית לכל אוזניה, מעבר למצב אוזניה אחת ללא ניתוק.
  \item \textbf{אודיו שידור Auracast}:Auracast. מכשיר אחד משדר ל-\textbf{אינספור}אזניות - שדה תעופה, מוזיאון, חדר כושר, כיתה.
  \item \textbf{מצב נוכחי}:2024-2026 חדשים תומכים, ישנים לא. חפש LE Audio/Auracast.
\end{itemize}

\end{tcbinfo}

\booksection{איך מתחבר}

: אם הרכב שלך תומך (רוב הרכבים החדשים בשנים האחרונות תומכים), חבר את הטלפון בכבל ל-USB ברכב, מסך הרכב יעבור אוטומטית לממשק CarPlay או Android Auto. חלק מהרכבים החדשים תומכים בחיבור אלחוטי, עוד יותר נוח.

\begin{itemize}
  \item \textbf{מה אפשר לעשות}: ניווט (הרבה יותר טוב מניווט הרכב), שיחות (ענייה מההגה), מוזיקה (אפליקציות המוזיקה בטלפון מוקרנות לרכב), שליחת WeChat קולית (Siri או Xiao Ai עוזרים לקרוא ולענות להודעות).
  \item \textbf{גם רכב ישן יכול}: אם הרכב שלך ישן ולא תומך, תמורת כמה מאות יואן קנה מסך CarPlay חיצוני, הדבק ללוח השעונים, חוויה דומה.
  \item \textbf{שבע, אבטחה ופרטיות בקישוריות}ככל שמחברים יותר מכשירים, שטח הסיכון לדליפת פרטיות גדל. כמה נקודות לתשומת לב:
\end{itemize}

\booksection{לפני הקרנה בציבור, בדוק קודם אם יש הודעות פרטיות בשורת ההתראות, אל תיתן לכל המשפחה לראות מה שלא צריך.}

לפני מסירת מכשיר למישהו אחר או מכירתו, בטל חיבור - בהגדרות הטלפון מצא את המכשיר, בחר “הסר” או “בטל חיבור”.

\begin{itemize}
  \item מצלמות בית חכם - קנה מותג עם הצפנה, אל תקנה חיקויים בעשרות יואן. שמעתי יותר מדי חדשות על מצלמות ששודרו על ידי האקרים.
  \item בוא ננסה: הקרן את מסך הטלפון לטלוויזיה, והגדר תרחיש חכם אחד
  \item ודא שהטלפון והטלוויזיה מחוברים לאותו Wi-Fi.
\end{itemize}

\begin{tcbpractice}{בטלוויזיה מצא את "שיקוף מסך" או "AirPlay" והפעל.}
\begin{enumerate}
  \item בטלפון: iPhone ממרכז הבקרה לוחץ על "שיקוף מסך"; אנדרואיד מושך למטה ומחפש "הקרנה אלחוטית".
  \item בחר תמונה שאתה הכי אוהב והקרן לטלוויזיה, תן למשפחה לראות.
  \item אם יש לך מכשירי בית חכם, נסה ליצור "מצב עזיבת הבית" באפליקציה: כפתור אחד לכיבוי אור, כיבוי מזגן, הפעלת רובוט השאיבה.
  \item פרק 18: גיבוי נתונים והעברה
  \item שכן שלי, הטלפון שלה נפל לבריכה וכשהוציאו אותו כבר לא נדלק. מה שהכאיב לה לא היה הטלפון עצמו, אלא התמונות שלא גובו שלוש שנים - כל הרישומים של הילד מלידה ועד שלמד ללכת. כששמעתי את זה הרגשתי רע מאוד, כי זה היה לגמרי מניע. בפרק זה נרכוש ברצינות פוליסת ביטוח לנתונים.
\end{enumerate}

\end{tcbpractice}

\begin{tcbinfo}{מקרה 1: אופטימיזציית חיי סוללה בפועל}
\begin{itemize}
  \item \textbf{תרחיש}: משתמש Xiaomi 14, חיי סוללה ירודים, מתרוקנת בחצי יום.
  \item \textbf{צעדים}:1) סוללה > דירוג; 2) App חריג; 3) 5G4G; 4) הגבל רקע; 5) מצב חיסכון; 6) הסר נרתיק.
  \item \textbf{תוצאה צפויה}:30-50\% יותר סוללה, נתוני בריאות.
\end{itemize}

\end{tcbinfo}
